\section*{Capítulo \ref{ch:g9:linfunc} --- Funciones lineales y afines}

\begin{solution}{exo:g9:linfunc:1}
Lineales: $f$ y $k$ (de la forma $ax$). Afines pero no lineales: $g$ (\omterm{prop:g9:linfunc:affinegraph}{pendiente} $3$,
ordenada en el origen $-5$) y $m$ (constante, \omterm{prop:g9:linfunc:affinegraph}{pendiente} $0$). Ninguna de las dos: $h$, por culpa
del cuadrado.
\end{solution}

\begin{solution}{exo:g9:linfunc:2}
$f(0) = -3$; $f(4) = 5$; $f(-1) = -5$. Y $f(x) = 9$ significa
$2x - 3 = 9$, luego $2x = 12$ y $x = 6$.
\end{solution}

\begin{solution}{exo:g9:linfunc:3}
$a = \frac{15}{6} = 2.5$, luego $f(x) = 2.5x$. Entonces $f(10) = 25$ y
$f(-2) = -5$.
\end{solution}

\begin{solution}{exo:g9:linfunc:4}
$f$ pasa por el origen con \omterm{prop:g9:linfunc:affinegraph}{pendiente} $2$; $g$ tiene la misma \omterm{prop:g9:linfunc:affinegraph}{pendiente} pero
corta el eje $y$ en $3$; $h$ decrece desde $(0, 4)$ con \omterm{prop:g9:linfunc:affinegraph}{pendiente}
$-1$. Las gráficas de $f$ y $g$ son \emph{\omterm{def:g6:lines:perp}{rectas paralelas}} (misma \omterm{prop:g9:linfunc:affinegraph}{pendiente},
distinta ordenada en el origen).
\end{solution}

\begin{solution}{exo:g9:linfunc:5}
Para $f$: \omterm{prop:g9:linfunc:affinegraph}{pendiente} $\frac{11 - 5}{3 - 1} = 3$, y después $5 = 3 \times 1 + b$
da $b = 2$: $f(x) = 3x + 2$.

Para $g$: \omterm{prop:g9:linfunc:affinegraph}{pendiente} $\frac{0 - 4}{2 - 0} = -2$, y $g(0) = 4$ ya es la
ordenada en el origen: $g(x) = -2x + 4$.
\end{solution}

\begin{solution}{exo:g9:linfunc:6}
Aumento del $25\%$: $120 \times 1.25 = 150$.

Descuento del $40\%$: $65 \times 0.6 = 39$.

Precio inicial $p$ con $p \times 0.75 = 69$:
$p = \frac{69}{0.75} = 92$.
\end{solution}

\begin{solution}{exo:g9:linfunc:7}
\emph{1.} $f(x) = 0.05x + 10$.

\emph{2.} La tarifa B sale más barata cuando $25 < 0.05x + 10$, es decir
$15 < 0.05x$, es decir $x > 300$. A partir de $301$ minutos ($5$ horas) al mes,
gana la tarifa B.
\end{solution}

\begin{solution}{exo:g9:linfunc:8}
\emph{1.} Al cabo de una hora: $2000 \times 1.1 = 2200$. Al cabo de dos horas:
$2200 \times 1.1 = 2420$.

\emph{2.} El segundo aumento se aplica a $2200$, no a $2000$: crecer
un $10\%$ dos veces multiplica por $1.1^2 = 1.21$, un aumento del $21\%$, no del
$20\%$.
\end{solution}

\begin{solution}{exo:g9:linfunc:9}
\emph{1.} \omterm{prop:g9:linfunc:affinegraph}{Pendiente}: $\frac{1 - 3}{6 - 2} = \frac{-2}{4} = -0.5$: la
función es decreciente.

\emph{2.} $f(x) = -0.5x + b$ con $f(2) = 3$: $-1 + b = 3$, luego $b = 4$
y $f(x) = -0.5x + 4$. Salida $0$: $-0.5x + 4 = 0$ da $x = 8$.
\end{solution}

\begin{solution}{exo:g9:linfunc:10}
\emph{1.} Las dos variaciones multiplican el precio por
\[
\left(1 + \frac{t}{100}\right)\left(1 - \frac{t}{100}\right)
= 1 - \frac{t^2}{10000}
\]
(tercera identidad con $a = 1$, $b = \frac{t}{100}$).

\emph{2.} Para $0 < t \leq 100$, $\frac{t^2}{10000} > 0$, así que el factor
es menor que $1$: el precio final es menor que el inicial. Para
$t = 10$: factor $1 - \frac{100}{10000} = 0.99$, es decir una bajada
total del $1\%$.
\end{solution}

\begin{solution}{pb:g9:linfunc:1}
\textbf{1.} \omterm{prop:g9:linfunc:affinegraph}{Pendiente}: $\frac{f(100) - f(0)}{100 - 0} =
\frac{212 - 32}{100} = 1.8$; ordenada en el origen $f(0) = 32$. Por tanto
\[
f(C) = 1.8\,C + 32 .
\]

\textbf{2.} $f(20) = 68\,^\circ$F (el famoso día templado);
$f(37) = 98.6\,^\circ$F; $f(-10) = 14\,^\circ$F.

\textbf{3.} De $F = 1.8C + 32$:
$C = \frac{F - 32}{1.8}$. Entonces $50\,^\circ$F
$\to \frac{18}{1.8} = 10\,^\circ$C, y $98.6\,^\circ$F
$\to \frac{66.6}{1.8} = 37\,^\circ$C.

\textbf{4.} $f(0) = 32 \neq 0$, y $f(20) = 68$ no es el
doble de $f(10) = 50$: no es lineal. La gráfica es una recta
que no pasa por el origen: $f$ es \emph{afín}, no lineal.

\textbf{5.} $1.8x + 32 = x$ da $0.8x = -32$, luego $x = -40$:
a cuarenta bajo \omterm{def:g1:counting:zero}{cero}, los dos termómetros marcan $-40$. En la gráfica, la
recta de $f$ corta ahí justamente a la diagonal $y = x$.

\textbf{6.} $p_A(k) = 1.5k + 3$ y $p_B(k) = 2k$. Para $4$ km:
$9$ euros frente a $8$: B más barato. Para $8$ km: $15$ frente a
$16$: A más barato.

\textbf{7.} $1.5k + 3 = 2k$ da $k = 6$: a seis kilómetros
los dos cuestan $12$ euros. Gráficamente, la recta de B (que pasa por el
origen y es más inclinada) empieza por debajo de la de A y la corta en
$(6, 12)$: B gana en los trayectos cortos y A en los largos.

\textbf{8.} $N = 7 \times 20 - 30 = 110$ cantos por minuto. De
$82 = 7T - 30$: $T = \frac{112}{7} = 16\,^\circ$C.

\textbf{9.} $v(t) = 600 - 15t$. Vale $240$:
$600 - 15t = 240$, $t = 24$ meses. Vale $0$: $t = 40$ meses.
Pasados los $40$ meses la fórmula se vuelve negativa --- pero el valor de un
teléfono se para en \omterm{def:g1:counting:zero}{cero}: todo modelo afín tiene un dominio en el que
tiene sentido.

\textbf{10.} Componiendo: $\times 1.25$ y después $\times 0.8$
multiplica por $1.25 \times 0.8 = 1$: el precio vuelve exactamente
a su punto de partida. No hay milagro: los porcentajes actúan sobre referencias
distintas (la bajada se aplica al precio ya subido) y aquí los
multiplicadores se cancelan por casualidad --- la historia general, con su
pérdida incorporada $1 - \frac{t^2}{10\,000}$, es el
\cref{exo:g9:linfunc:10}.

\textbf{11.} \omterm{prop:g9:linfunc:affinegraph}{Pendiente}: $\frac{17 - 20}{30} = -0.1$ cm por minuto;
$h(t) = 20 - 0.1t$. Se apaga en $h(t) = 0$: $t = 200$ minutos
--- tres horas y veinte minutos de luz.

\textbf{12.} El segundo salió $10$ km por delante (ordenada en el origen); el
primero es más rápido ($24 > 18$ km/h). Alcance:
$24t = 10 + 18t$ da $6t = 10$, $t = \frac53$ h $= 1$ h $40$
min, en el kilómetro $24 \times \frac53 = 40$.

\textbf{13.} $f(x) + g(x) = (2 - 0.5)x + (3 + 7) = 1.5x + 10$:
afín, con \omterm{prop:g9:linfunc:affinegraph}{pendiente} la \emph{\omterm{def:g1:addition:def}{suma} de las \omterm{prop:g9:linfunc:affinegraph}{pendientes}} y ordenada en el origen
la \omterm{def:g1:addition:def}{suma} de las ordenadas en el origen --- se ve al agrupar los términos semejantes
en $(mx + p) + (m'x + p')$.

\textbf{14.} Si $m \neq 0$: exactamente una solución,
$x = -\frac pm$ --- una recta no horizontal corta el eje
una vez. Si $m = 0$, $p \neq 0$: no hay solución --- una recta horizontal
por encima o por debajo del eje nunca lo toca. Si $m = 0$, $p = 0$:
todo $x$ es solución --- la recta \emph{es} el eje.

\textbf{15.} Primer terna: tasas $\frac{9 - 5}{3 - 1} = 2$ y
$\frac{17 - 9}{7 - 3} = 2$: iguales, los puntos están alineados (sobre
$y = 2x + 3$). Segunda terna: $\frac{9-5}{2} = 2$, pero
$\frac{16 - 9}{3} = \frac73 \neq 2$: no están alineados. Las \omterm{def:g9:linfunc:affine}{funciones
afines} son exactamente las de tasa de variación constante ---
un solo número, la \omterm{prop:g9:linfunc:affinegraph}{pendiente}, cuenta toda la historia; en las curvas, la
tasa cambia de un punto a otro, y perseguirla lleva
a la derivada, en el \omterm{def:g6:measure:volume}{volumen} de secundaria superior.
\end{solution}
